\documentclass[11pt]{article}
\usepackage[utf8]{inputenc}
\usepackage{amsmath,amssymb,amsthm}
\usepackage{hyperref}
\usepackage{listings}
\usepackage{xcolor}
\usepackage[margin=1in]{geometry}

\lstset{
    basicstyle=\ttfamily\small,
    breaklines=true,
    frame=single,
    backgroundcolor=\color{gray!5},
    numbers=none,
    xleftmargin=4pt,
    xrightmargin=4pt,
    aboveskip=8pt,
    belowskip=8pt
}

\newtheorem{conjecture}{Conjecture}

\title{Empirical Investigation of an Open Conjecture:\\
A unitary divisor d of a positive integer n is a divisor such that d and n/d are}
\author{Agentic NL$\to$Lean~4 Pipeline \\ \small Job \#23}
\date{\today}

\begin{document}
\maketitle

\begin{abstract}
This report documents the empirical investigation of an open mathematical conjecture
that could not be formally proved or disproved in Lean~4 with Mathlib.
Numerical experiments were conducted to gather evidence for or against the conjecture.
The empirical verdict is: \textbf{\textcolor{green!70!black}{Empirically Supported}}.
The conjecture remains formally open.
\end{abstract}

\section{Conjecture Statement}

\begin{conjecture}
\begin{lstlisting}
A unitary divisor d of a positive integer n is a divisor such that d and n/d are coprime. A positive integer n is called a unitary perfect number if it is the sum of its unitary divisors (aside from n itself). Conjecture: there are only finitely many unitary perfect numbers.
\end{lstlisting}
\end{conjecture}

\section{Status}

\begin{center}
\fbox{\parbox{0.8\textwidth}{%
\textbf{Formal Status:} OPEN --- no Lean~4 proof or disproof was found.\\[4pt]
\textbf{Empirical Verdict:} \textcolor{green!70!black}{Empirically Supported}
}}
\end{center}

The pipeline attempted formal verification in Lean~4 with Mathlib but was unable to
produce a compiling proof or disproof. Empirical testing was then conducted to gather
numerical evidence.

\section{Basic Empirical Testing}

The following output was produced by the basic numerical experiment:

\begin{lstlisting}
=== EXPERIMENT PLAN ===

Conjecture: There are only finitely many unitary perfect numbers (UPNs).

A unitary divisor d of n satisfies gcd(d, n/d) = 1.
Equivalently, if n = prod p_i^{a_i}, the unitary divisors are
products of 'full' prime-power blocks p_i^{a_i} (each block is
either fully included or not). So sigma*(n) = prod (1 + p_i^{a_i}).
n is a unitary perfect number iff sigma*(n) = 2n,
i.e., sum of proper unitary divisors = n.

Known UPNs (historical, Wall 1975 etc.):
  6, 60, 90, 87360, 146361946186458562560000
No sixth UPN has been found despite extensive search.

Our experiment has FIVE complementary angles:

(1) Exhaustive search up to a large bound N_max (integer-by-integer),
    using a sieve computing sigma*(n) efficiently, confirming the
    first four known UPNs and searching for any new small UPN.

(2) Structured search over highly-composite-like candidates:
    numbers of the form 2^a * 3^b * 5^c * 7^d * ... with random
    exponents, up to 1e22. This samples a huge combinatorial
    space near the "abundancy frontier", where UPNs must live
    (sigma*(n)/n = 2 exactly).

(3) Distribution of the unitary-abundancy ratio u(n) = sigma*(n)/n.
    Plot histograms and show how rarely u(n) lands exactly on 2.
    This gives a heuristic density argument: the chance of
    u(n) = 2 decreases as n grows (abundancy values become
    arithmetically constrained), supporting finiteness.

(4) 2-adic obstruction check: a classical result shows that
    every UPN must be even and that the exponent of 2 matters.
    We verify that the known UPNs satisfy 2||n-like constraints
    (2^a || n for specific a) and examine which (a, odd part)
    configurations can possibly give sigma*(n) = 2n.

(5) Heuristic growth model: consider the count U(x) of UPNs <= x.
    We empirically estimate log U(x) vs log x from the known
    values; if U(x) grows sub-logarithmically or is bounded,
    finiteness is supported.

Verdict rule: since the conjecture is an open infinity/finiteness
statement, empirical evidence can only SUPPORT or weaken it.
We declare EMPIRICALLY SUPPORTED if (a) no new UPN is found
in extensive search, (b) density of UPNs is vanishingly small
and decreasing, (c) known UPNs fit a highly sparse pattern.

(1) Exhaustive sieve search up to N_max = 200000
    UPNs found up to 200000: [6, 60, 90, 87360]
    (Expected from literature: 6, 60, 90, 87360)

(2) Structured search over smooth (highly composite) candidates
    Trials: 200000, UPNs found in structured search: 1
    Distinct UPNs found: [6] ...
    Verification: sigma*(known 5th UPN)/n = 2.0

(3) Distribution of u(n) = sigma*(n)/n
    Small-n sample size: 199999
    mean u(n) = 1.3684
    max  u(n) = 3.2224
    P(u(n) >= 2) = 7.04604e-02
    P(u(n) == 2) (exact UPN density) = 2.00001e-05
    density in [2,100) = 3.061e-02 (3 UPNs)
    density in [100,1000) = 0.000e+00 (0 UPNs)
    density in [1000,10000) = 0.000e+00 (0 UPNs)
    density in [10000,100000) = 1.111e-05 (1 UPNs)
    density in [100000,200000) = 0.000e+00 (0 UPNs)

(4) Structural/2-adic pattern of known UPNs
    n=6: 2-exponent=1, odd part=3, omega(n)=2
    n=60: 2-exponent=2, odd part=15, omega(n)=3
    n=90: 2-exponent=1, odd part=45, omega(n)=3
    n=87360: 2-exponent=6, odd part=1365, omega(n)=5
    n=146361946186458562560000: 2-exponent=18, odd part=558326515909036875, omega(n)=12
    All known/found UPNs even: True

(5) Growth of U(x) = #{UPNs <= x}
    log(x) of known UPNs: [ 1.79175947  4.09434456  4.49980967 11.37779279 53.34036959]
    U(x) increments: [1 2 3 4 5]
    gaps in log(x) between successive UPNs: [ 2.30258509  0.40546511  6.87798312 41.9625768 ]
    Very large and GROWING log-gaps -> strong sparsity signal.

=== SUMMARY ===
Runtime: 3.33 s
(1) Exhaustive search up to 200000: found [6, 60, 90, 87360] (matches known: 6, 60, 90, 87360)
(2) Structured search: 200000 trials over smooth numbers up to 1e22, NEW UPNs found = []
(3) Empirical P(u(n) = 2) over [2, 200000] = 2.000e-05 (only 4 hits in 2
... [truncated]
\end{lstlisting}


\section{Experiment Code (Basic)}

\begin{lstlisting}[language=Python]
import matplotlib
matplotlib.use("Agg")
import matplotlib.pyplot as plt
import numpy as np
import math
import random
from sympy import factorint, divisors
from collections import Counter
import time

print("=== EXPERIMENT PLAN ===")
print("""
Conjecture: There are only finitely many unitary perfect numbers (UPNs).

A unitary divisor d of n satisfies gcd(d, n/d) = 1.
Equivalently, if n = prod p_i^{a_i}, the unitary divisors are
products of 'full' prime-power blocks p_i^{a_i} (each block is
either fully included or not). So sigma*(n) = prod (1 + p_i^{a_i}).
n is a unitary perfect number iff sigma*(n) = 2n,
i.e., sum of proper unitary divisors = n.

Known UPNs (historical, Wall 1975 etc.):
  6, 60, 90, 87360, 146361946186458562560000
No sixth UPN has been found despite extensive search.

Our experiment has FIVE complementary angles:

(1) Exhaustive search up to a large bound N_max (integer-by-integer),
    using a sieve computing sigma*(n) efficiently, confirming the
    first four known UPNs and searching for any new small UPN.

(2) Structured search over highly-composite-like candidates:
    numbers of the form 2^a * 3^b * 5^c * 7^d * ... with random
    exponents, up to 1e22. This samples a huge combinatorial
    space near the "abundancy frontier", where UPNs must live
    (sigma*(n)/n = 2 exactly).

(3) Distribution of the unitary-abundancy ratio u(n) = sigma*(n)/n.
    Plot histograms and show how rarely u(n) lands exactly on 2.
    This gives a heuristic density argument: the chance of
    u(n) = 2 decreases as n grows (abundancy values become
    arithmetically constrained), supporting finiteness.

(4) 2-adic obstruction check: a classical result shows that
    every UPN must be even and that the exponent of 2 matters.
    We verify that the known UPNs satisfy 2||n-like constraints
    (2^a || n for specific a) and examine which (a, odd part)
    configurations can possibly give sigma*(n) = 2n.

(5) Heuristic growth model: consider the count U(x) of UPNs <= x.
    We empirically estimate log U(x) vs log x from the known
    values; if U(x) grows sub-logarithmically or is bounded,
    finiteness is supported.

Verdict rule: since the conjecture is an open infinity/finiteness
statement, empirical evidence can only SUPPORT or weaken it.
We declare EMPIRICALLY SUPPORTED if (a) no new UPN is found
in extensive search, (b) density of UPNs is vanishingly small
and decreasing, (c) known UPNs fit a highly sparse pattern.
""")

t0 = time.time()

# ---------- (1) Sieve-based exhaustive search ----------
# sigma*(n) = product over prime powers p^a || n of (1 + p^a)
# Use a sieve: for each prime power p^a with p^a <= N, multiply
# sigma*[m] by (1 + p^a) for every m that has p^a exactly as its
# p-part. Implement by iterating primes and exponents.

def unitary_sigma_sieve(N):
    sig = np.ones(N + 1, dtype=np.int64)
    sig[0] = 0
    # sieve smallest prime factor
    spf = np.zeros(N + 1, dtype=np.int64)
    for i in range(2, N + 1):
        if spf[i] == 0:
            for j in range(i, N + 1, i):
                if spf[j] == 0:
                    spf[j] = i
    # For each n, factor via spf and compute sigma*
    for n in range(2, N + 1):
        m = n
        s = 1
        while m > 1:
            p = spf[m]
            pa = 1
            while m % p == 0:
                m //= p
                pa *= p
            s *= (1 + pa)
        sig[n] = s
    return sig

N_max = 200000
print(f"(1) Exhaustive sieve search up to N_max = {N_max}")
sig_star = unitary_sigma_sieve(N_max)
upn_small = [n for n in range(2, N_max + 1) if sig_star[n] == 2 * n]
print(f"    UPNs found up to {N_max}: {upn_small}")
print(f"    (Expected from literature: 6, 60, 90, 87360)")

# ---------- (2) Structured search over smooth numbers ----------
print("\n(2) Structured search over smooth (highly composite) candidates")

def sigma_star_from_factorization(fac):
    s = 1
    for p, a in fac.items():
        s *= (1 + p ** a)
    return s

def n_from_factorization(fac):
    n = 1
    for p, a in fac.items():
        n *= p ** a
    return n

primes_small = [2, 3, 5, 7, 11, 13, 17, 19, 23, 29, 31, 37, 41, 43, 47]
max_exp_per_prime = {2:30, 3:20, 5:15, 7:12, 11:10, 13:9, 17:8, 19:8,
                     23:7, 29:7, 31:7, 37:6, 41:6, 43:6, 47:6}
BOUND = 10**22

rng = random.Random(42)
trials = 200000
upn_structured = []
ratios = []
for _ in range(trials):
    k = rng.randint(2, 10)
    ps = rng.sample(primes_small, k)
    fac = {}
    n = 1
    for p in ps:
        a = rng.randint(1, max_exp_per_prime[p])
        if n * p ** a > BOUND:
            a = 1
            if n * p > BOUND:
                continue
        fac[p] = a
        n *= p ** a
    if n < 2:
        continue
    s = sigma_star_from_factorization(fac)
    r = s / n
    ratios.append(r)
    if s == 2 * n:
        upn_structured.append(n)

upn_structured = sorted(set(upn_structured))
print(f"    Trials: {trials}, UPNs found in structured search: {len(upn_structured)}")
if upn_structu
# ... [truncated]
\end{lstlisting}


\section{Conclusion}

The conjecture remains formally open. Numerical experiments \textbf{support} the conjecture --- no counterexamples were found across all tested parameter ranges.
Further investigation (both formal and empirical) is warranted.

\end{document}
